\section{Sustainable Development Goal Alignment}

WealthGenie is designed as an educational platform aligned with four United Nations Sustainable Development Goals. The platform does not measure, achieve or guarantee SDG targets; rather, its features are designed to support the objectives of each goal through financial education and transparent decision support.

\subsection{SDG 4: Quality Education}

SDG 4 targets inclusive and equitable quality education that promotes lifelong learning opportunities for all. The Learn module is WealthGenie's primary contribution to this goal. The module delivers 23 structured financial concepts organized into five branches: Money Basics, Saving, Growing Money, Investing and Financial Safety. Each concept is supported by a versioned lesson file containing objectives, explanations, worked examples, key takeaways, a three-question quiz with explanations, citations and a review record~\cite{learning_content}. Lessons are served from the backend through the curated lesson loader without requiring an AI provider, ensuring that reading published content, examples, quizzes and citations never depends on an external language model.

The curriculum covers income, expenses and budgeting; saving money and emergency funds; inflation; simple and compound interest; risk and return; bank deposits, fixed and recurring deposits, stocks, mutual funds, index funds, SIPs, diversification, fees and expense ratios, investment time horizons and risk tolerance; UPI safety, digital financial fraud, borrowing and interest, and credit scores. All content is grounded in verified publications of the RBI FAME, SEBI investor education documents and AMFI disclosures. Progress persistence through the learning progress API enables learners to track completion across sessions, supporting the lifelong learning dimension of SDG 4.

\subsection{SDG 8: Decent Work and Economic Growth}

SDG 8 promotes sustained, inclusive and sustainable economic growth, full and productive employment and decent work for all. Informed personal financial decision-making contributes to household economic stability, which in turn supports broader economic participation. The Assessment module helps users understand their financial position by computing monthly surplus, savings rate, emergency fund coverage and goal projections from their own reported income, expenses, savings and obligations. By making these indicators explicit and traceable, the platform helps users recognize the relationship between current financial behaviour and future outcomes, supporting more deliberate financial planning that is consistent with sustained economic well-being.

\subsection{SDG 10: Reduced Inequalities}

SDG 10 aims to reduce inequality within and among countries. Financial exclusion is a significant dimension of economic inequality, and barriers to understanding formal financial products disproportionately affect lower-income and first-time investors. WealthGenie is designed to be accessible: it requires no minimum investable amount, provides clear plain-language explanations alongside technical metrics, and exposes the reasoning behind every recommendation. The platform's Learn and Assessment modules are available without authentication, and the recommendation engine degrades gracefully when material inputs are missing, reporting what is unknown rather than fabricating values. By reducing the knowledge barriers that separate informed participants from the formal financial system, the platform aligns with the inclusion objectives of SDG 10.

\subsection{SDG 12: Responsible Consumption and Production}

SDG 12 calls for sustainable consumption and production patterns. In personal finance, responsible consumption translates into awareness of costs, risk and the long-term consequences of investment choices. The recommendation engine operationalizes this principle by design: historical returns are displayed for transparency but are explicitly excluded from the scoring function, which instead rewards risk alignment, low expense ratios, contained volatility and drawdown, return consistency and data completeness. The paper trading module provides a zero-risk environment in which users can experiment with buying and selling before committing real capital. The Digital Twin enables users to explore the effects of changing financial assumptions on projected outcomes. Together, these features encourage considered, evidence-based financial behaviour rather than speculative decision-making, consistent with the responsible consumption objectives of SDG 12.